\section*{The whole project on one page}
\addcontentsline{toc}{section}{The whole project on one page}
\markboth{The whole project on one page}{The whole project on one page}

\textbf{The question.} A neural network's weights and the values it computes while running
(activations) sit in memory. Sometimes a single bit in memory flips by accident. In an ordinary
number format, a flipped bit damages one number. In \emph{microscaling} (MX) formats, a block of
32 numbers shares one 8-bit scale, so a flip in that scale damages all 32 numbers at once. How much
do these flips actually hurt a network's predictions, and which formats, bits, layers and tensors
are most dangerous?

\textbf{What you built.} A fault-injection tool called \code{mxfi}. It encodes weights and activations
into real MX bit patterns, flips a chosen bit in an element or in a scale, decodes, runs the network,
and compares the predictions with the fault-free ones.

\textbf{What you ran.} About 60 trained networks (ResNet8, RepVGG-A0, a small Vision Transformer;
CIFAR-10 and CIFAR-100), five element formats, block sizes 8 to 64, weights and activations.
Roughly 1.44 million sampled weight faults in 119 campaigns, two \emph{exhaustive} campaigns that
flip every single stored bit of a ResNet8 (1.12 million injections), 390{,}000 multi-bit injections,
and further campaigns for activations, a NaN guard and CIFAR-100. The work ran on your laptop and on a
two-A100 server at IIT Kanpur.

\textbf{What you found.}
\begin{enumerate}
\item Formats with the same accuracy differ in fault vulnerability by up to three orders of magnitude.
 Measured per inference, e3m2 $<$ e4m3 $<$ e5m2 on every network tested.
\item The main cause is the NaN/Infinity codes of the 8-bit formats. On the larger networks 97 to 100\%
 of element damage comes from flips that land on such a code. Decoding those codes as zero on read cuts
 element damage by at least $29\times$ on RepVGG and the transformer.
\item A fault in the shared scale does $13\times$ to over $1000\times$ the damage of a fault in an element.
\item The usual ``did any test image change?'' metric mostly measures whether the test set contains a
 borderline image. You replaced it with a per-inference rate.
\item The severity score in your proposal does not predict failures. A gradient-based score does, and sampling
 with it needs $57$ to $120\times$ fewer injections.
\item Both open design choices in the proposal are settled: single-bit injection is enough, and exact
 enumeration of the code table beats TreeFI-style learned value intervals.
\end{enumerate}

\textbf{What you corrected along the way.} Eleven earlier conclusions were revised or withdrawn,
for reasons that come down to three causes: a misleading metric, too few trained networks, and
mismatched checkpoints and a codec bug. Section~\ref{sec:mistakes} lists all of them.

\textbf{What was not done.} ImageNet and DeiT were not reachable from any machine you could use, so
CIFAR-100 stands in for them. The optional small language model and MXINT8 were not studied.

\medskip
\begin{center}
\small
\begin{tabular}{@{}lr@{}}
\toprule
Report & 13 pages, \code{docs/report.pdf} \\
Findings log & 72 numbered findings, \code{docs/findings.md} \\
Library & \code{mxfi/}: 11 modules plus \code{__init__}, about 2{,}300 lines \\
Experiment scripts & \code{experiments/}: 29 scripts, about 5{,}200 lines \\
Tests & 4 files, about 1{,}100 lines \\
Repository & \code{github.com/peeyushagarwal2004/CS496_UGP} \\
\bottomrule
\end{tabular}
\end{center}
\src{README.md, docs/report.tex, docs/findings.md; line counts measured with wc on the repository.}

\newpage
